\chapter{Resoluções das atividades propostas}
\label{cap:solucoes}

\setcounter{equation}{0}
\renewcommand{\theequation}{S.\arabic{equation}}
\newcounter{solucaolistingbase}
\setcounter{solucaolistingbase}{\value{lstlisting}}
\renewcommand{\thelstlisting}{S.\number\numexpr
  \value{lstlisting}-\value{solucaolistingbase}\relax}

As resoluções são organizadas por capítulo e mantêm a numeração dos exercícios no texto principal. Em cada solução, o desenvolvimento matemático precede a substituição numérica. O Python é usado como verificação quando acrescenta valor ao raciocínio.

\section*{Capítulo 1 — Ferramentas matemáticas para a Física}
\addcontentsline{toc}{section}{Capítulo 1 — Ferramentas matemáticas para a Física}

\subsection*{Exercício 1}

Para escrever o diâmetro em notação científica, deslocamos a vírgula sete posições para a direita:

\begin{equation}
  \SI{0.00000032}{\meter}
  =\SI{3.2e-7}{\meter}.
\end{equation}

Como \(\SI{1}{\meter}=\SI{1e9}{\nano\meter}\), multiplicamos o valor em metros por \(10^9\):

\begin{equation}
  (3{,}2\times10^{-7}\,\si{\meter})
  \left(10^9\,\frac{\si{\nano\meter}}{\si{\meter}}\right)
  =3{,}2\times10^2\,\si{\nano\meter}
  =\SI{320}{\nano\meter}.
\end{equation}

Portanto, o diâmetro é \(\boxed{\SI{3.2e-7}{\meter}=\SI{320}{\nano\meter}}\).

\subsection*{Exercício 2}

Se a energia \(E\) é proporcional a \(q^2\), podemos escrever \(E=Cq^2\), onde \(C\) permanece constante. Ao substituir \(q\) por \(3q\), obtemos

\begin{equation}
  E'=C(3q)^2=9Cq^2=9E.
\end{equation}

Logo, a energia é \(\boxed{9}\) vezes maior.

\subsection*{Exercício 3}

Para uma intensidade \(I\) inversamente proporcional a \(r^2\), temos \(I=C/r^2\). Assim,

\begin{equation}
  \frac{I(\SI{2}{\meter})}{I(\SI{6}{\meter})}
  =\frac{C/2^2}{C/6^2}
  =\frac{6^2}{2^2}
  =9.
\end{equation}

A intensidade a \(\SI{2}{\meter}\) é, portanto, \(\boxed{9}\) vezes a intensidade a \(\SI{6}{\meter}\). Equivalentemente, \(I(\SI{6}{\meter})=I(\SI{2}{\meter})/9\).

\subsection*{Exercício 4}

Tomando o ângulo a partir do eixo horizontal,

\begin{align}
  F_x &= F\cos\theta
       =\SI{120}{\newton}\cos 40^\circ
       \approx\SI{91.9}{\newton},\\
  F_y &= F\sin\theta
       =\SI{120}{\newton}\sin 40^\circ
       \approx\SI{77.1}{\newton}.
\end{align}

As duas componentes são positivas porque o vetor está no primeiro quadrante. Portanto,

\begin{equation}
  \boxed{\vetor{F}\approx
  (\SI{91.9}{\newton})\,\hat{\vetor{\imath}}
  +(\SI{77.1}{\newton})\,\hat{\vetor{\jmath}}}.
\end{equation}

Como verificação, \(\sqrt{F_x^2+F_y^2}\approx\SI{120}{\newton}\).

\subsection*{Exercício 5}

Para converter graus em radianos, multiplicamos por \(\pi/180^\circ\):

\begin{equation}
  150^\circ\frac{\pi\,\si{\radian}}{180^\circ}
  =\boxed{\frac{5\pi}{6}\,\si{\radian}}.
\end{equation}

No sentido inverso,

\begin{equation}
  \frac{5\pi}{6}\,\si{\radian}
  \frac{180^\circ}{\pi\,\si{\radian}}
  =\boxed{150^\circ}.
\end{equation}

Os resultados são consistentes porque as duas expressões representam o mesmo ângulo.

\subsection*{Exercício 6}

A velocidade é a derivada da posição em relação ao tempo:

\begin{equation}
  v(t)=\frac{\dd x}{\dd t}=2-6t.
\end{equation}

A aceleração é a derivada da velocidade:

\begin{equation}
  a(t)=\frac{\dd v}{\dd t}=-6.
\end{equation}

Se \(x\) é medido em metros e \(t\) em segundos, cada termo de \(x(t)\) precisa ter unidade de comprimento. Portanto, o termo constante \(4\) tem unidade de metro, o coeficiente \(2\) tem unidade de \(\si{\meter\per\second}\), e o coeficiente \(-3\) tem unidade de \(\si{\meter\per\second\squared}\). Assim,

\begin{align}
  v(t)&=\boxed{(2-6t)\,\si{\meter\per\second}},\\
  a(t)&=\boxed{-6\,\si{\meter\per\second\squared}}.
\end{align}

Na primeira expressão, os números carregam as unidades necessárias para que ambos os termos representem velocidade.

\subsection*{Exercício 7}

O deslocamento é a integral da velocidade:

\begin{align}
  \Delta x
  &=\int_0^4(5-2t)\,\dd t\\
  &=\left[5t-t^2\right]_0^4\\
  &=20-16\\
  &=\boxed{\SI{4}{\meter}}.
\end{align}

No gráfico \(v\times t\), a velocidade se anula em \(t=\SI{2.5}{\second}\). Entre \(0\) e \(\SI{2.5}{\second}\), a área triangular positiva vale

\begin{equation}
  A_+=\frac{1}{2}(\SI{2.5}{\second})(\SI{5}{\meter\per\second})
  =\SI{6.25}{\meter}.
\end{equation}

Entre \(\SI{2.5}{\second}\) e \(\SI{4}{\second}\), a região fica abaixo do eixo e contribui com

\begin{equation}
  A_-=-\frac{1}{2}(\SI{1.5}{\second})(\SI{3}{\meter\per\second})
  =-\SI{2.25}{\meter}.
\end{equation}

Logo, \(\Delta x=A_++A_-=\SI{4}{\meter}\). O sinal positivo indica que, apesar da inversão do movimento após \(\SI{2.5}{\second}\), o deslocamento no sentido positivo foi maior.

\subsection*{Exercício 8}

Os pontos críticos de uma função diferenciável ocorrem onde \(f'(x)=0\). Para

\begin{equation}
  f(x)=x^3-6x^2+9x+2,
\end{equation}

temos

\begin{equation}
  f'(x)=3x^2-12x+9=3(x-1)(x-3).
\end{equation}

Portanto, os valores críticos são \(x=1\) e \(x=3\). A segunda derivada é

\begin{equation}
  f''(x)=6x-12.
\end{equation}

Em \(x=1\), \(f''(1)=-6<0\), logo há um máximo local. Como \(f(1)=6\), o ponto é \(\boxed{(1,6)}\).

Em \(x=3\), \(f''(3)=6>0\), logo há um mínimo local. Como \(f(3)=2\), o ponto é \(\boxed{(3,2)}\).

\subsection*{Exercício 9}

Definimos o excesso de temperatura em relação ao ambiente como

\begin{equation}
  \Theta(t)=T(t)-T_{\mathrm{amb}}
  =\Delta T_0\e^{-t/\tau}.
\end{equation}

Derivando em relação ao tempo,

\begin{align}
  \frac{\dd\Theta}{\dd t}
  &=-\frac{\Delta T_0}{\tau}\e^{-t/\tau}\\
  &=-\frac{1}{\tau}\Theta.
\end{align}

Como \(\Theta=T-T_{\mathrm{amb}}\), concluímos que

\begin{equation}
  \boxed{\frac{\dd}{\dd t}(T-T_{\mathrm{amb}})
  =-\frac{1}{\tau}(T-T_{\mathrm{amb}})}.
\end{equation}

A taxa é proporcional ao próprio excesso de temperatura. O sinal negativo indica que, para um corpo inicialmente mais quente que o ambiente, o excesso diminui com o tempo.

\subsection*{Exercício 10}

A inclinação da reta posição-tempo é

\begin{equation}
  v=\frac{\Delta x}{\Delta t}
  =\frac{\SI{25}{\meter}-\SI{7}{\meter}}
         {\SI{8}{\second}-\SI{2}{\second}}
  =\boxed{\SI{3}{\meter\per\second}}.
\end{equation}

Essa inclinação representa a velocidade constante. Escrevendo a posição como \(x(t)=x_0+vt\) e usando o ponto \((\SI{2}{\second},\SI{7}{\meter})\), obtemos

\begin{equation}
  \SI{7}{\meter}=x_0+
  (\SI{3}{\meter\per\second})(\SI{2}{\second}),
\end{equation}

de onde

\begin{equation}
  \boxed{x_0=\SI{1}{\meter}}.
\end{equation}

Portanto, a equação da posição é \(\boxed{x(t)=1+3t}\), com \(x\) em metros e \(t\) em segundos.

\subsection*{Verificação numérica em Python}

O programa a seguir confere os resultados numéricos dos exercícios 1, 4, 5, 7, 8 e 10. Ele não substitui os desenvolvimentos anteriores: apenas executa as operações depois que o modelo já foi estabelecido.

\lstinputlisting[
  caption={Verificação numérica de resultados selecionados do Capítulo 1.},
  label={cod:verificacao-solucoes-cap01}
]{codigos/solucoes/verificacao_cap01.py}

A saída esperada é

\lstinputlisting[
  style=saidaesperada,
  caption={Saída esperada da verificação das soluções do Capítulo 1.}
]{codigos/solucoes/saidas/verificacao_cap01.txt}
